\documentclass{article} %210 mm × 297 mm
\usepackage[utf8]{inputenc}
\usepackage[a4paper,left=1.5cm, right=1.5cm, top=2cm, bottom=2cm]{geometry}
\usepackage{graphicx}
%\usepackage{blindtext}
\usepackage{multirow}
\usepackage{mhchem}
\usepackage{url}
\usepackage{subfigure}
\usepackage{amsfonts,latexsym} % para tener disponibilidad de diversos simbolos
\usepackage{enumerate}
%\usepackage{booktabs}
\usepackage{float}
%\usepackage{threeparttable}
\usepackage{array,colortbl}
%\usepackage{ifpdf}
%\usepackage{rotating}
\usepackage{cite}
\usepackage{stfloats}
\usepackage{url}
\usepackage{listings}
\usepackage{fancyhdr}

\usepackage{color}
%\usepackage{hyperref}
%\usepackage{wrapfig}
\usepackage{array}
%\usepackage{multirow}
%\usepackage{adjustbox}
%\usepackage{nccmath}
\usepackage[dvipsnames]{xcolor}


\newcommand{\titulo}{Präsenzübungsblatt 5; Statistik}
\newcommand{\nombre}{Lil'Jana}
\newcommand{\FCE}{\textbf{SoSe 2025}}
\newcommand{\dat}{\textbf{Tut 13.05.2025}}
\newcommand{\p}[0]{\mathbb{P}}
 

\usepackage{fancyhdr}
\pagestyle{fancy}
\fancyhead{}
\fancyfoot{}
\fancyhead[RO]{\small\dat}
\fancyhead[LO]{\small\FCE}
\fancyfoot[RO,LE] {\thepage}

\setcounter{page}{1} %Número de la página, la editorial lo cambiará



\usepackage[onehalfspacing]{setspace}
\begin{document}


\begin{tabular}{p{12cm}}
{{\Large\bf\titulo}}\\
\vspace{0.2cm}\\
 {\large\nombre}\\
\end{tabular}
\vspace{0.2cm}\\
\section{Aufgabe 5.1}
\[
X(\omega) = 
\begin{cases}
0, \quad wenn\  \omega  =B \\
1, \quad wenn\  \omega  =B\ oder\ \omega = C
\end{cases}, \quad 
Y(\omega) = 
\begin{cases}
1, \quad wenn\  \omega  =B \\
0, \quad wenn\  \omega  =B\ oder\ \omega = C
\end{cases}
\]
\subsection*{(a)}
Man muss benutzen, dass es gleichverteilt ist, also $\mathbb{P}(X=0) = \frac{1}{3}$ 
\begin{table}[H]
    \centering
   \begin{tabular}{c|cc}
   $\mu_{x,y}$ & X= 0 & X=1 \\
    Y = 0 & 0 & $\frac{2}{3}$\\
    Y =1 & $\frac{1}{3}$ & 0
\end{tabular}
\end{table}
\caption{Plausibilitätscheck: Wenn wir alle Zellen summieren, sollte 1 raus kommen }
\begin{align*}
    \p(X=0, Y=0) &  = \p(\underbrace{X^{-1} (0)}_{=\{\omega \in \Omega: X \in (\omega) = 0\}} \cap Y ^{-1} (0))   \\
    & = \p (\{A\} \cap \{B,C\}) \\
    & = \p(\emptyset) = 0 \\ \\
    \p(X=1, Y = 0) & = \p (X^{-1}(1) \cap Y^{-1}(0) \\
    & = \p(\{B,C\} \cap \{B,C\})  \\
    & = \p(\{B,C\}) \\
    & = \mu(B) + \mu (C) = \frac{2}{3}
\end{align*}
\subsection*{(b)}
Randverteilung $\mu_X$ = Verteilung von $X$
\[
\mu_X : \{0,1\} \rightarrow [0,1]
\]
\[
\Omega \overset{X}{\rightarrow} \{0,1\} \overset{\mu_X}{\rightarrow} [0,1]
\]
\begin{align*}
    \mu_X(0) & = \frac{1}{3} \\
    & = \p(x=0) \\
    & = \mu_X(0) = \mu_{X,Y} ((0,0)) + \mu_{X,Y}((0,1)) \\
    & = 0+\frac{1}{3} \\
    & =\p (X=0, X=0) + \p (X=0, Y=1) \\
    & = \sum_{\substack{\omega \in \{0,1\}^2 \\ \omega_1 = 0}}  \mu_{X,Y}(\omega) \\ \\
\end{align*}
\begin{align*}
      \mu_X(1) & = \frac{2}{3} + 0 \\
    & = \frac{2}{3} \\ \\
    \mu_X(0) & = \sum_{ \substack{ \omega \in \{A,B,C\}\\ X\omega = 0}}  \mu
\end{align*}
\subsection*{(c)}
Verteilungfunktion von $X$. 
\[
F_x: \mathbb{R} \rightarrow [0,1]
\] 
\[
F_x(s) = \{x \leq s \} = \begin{cases}
0  & , s < 0 \\
\p(X=0) = \frac{1}{3} & , 0 \leq s < 1 \\
\p(x\leq1)^*=1& , 1 \leq s
\end{cases}
\]
\[
* \p(x\leq1) = \p(x=0\ oder\ x =1)
\]

\section*{Besprechung 3.1}
b) 
\[

\]
\end{document}